One question we weigh differently among ourselves, and instead of smoothing the difference over in the body, the afterword gives it the treatment the hard cases get, for it is a live question about the shape of the Church, it excites us, and the reader deserves the record.

The scriptural record first. The New Testament uses bishop and presbyter for the same men. St. Paul summons the presbyters of Ephesus and tells them the Holy Spirit has made them overseers of the flock, episkopoi (Acts 20:17, 28). He greets the bishops and deacons of a single town (Philippians 1:1). He tells Titus to appoint presbyters and gives the qualifications of a bishop in the same breath (Titus 1:5–7). St. Peter writes as a fellow elder to the elders (1 Peter 5:1), and St. John twice opens his letters as the presbyter. Lightfoot, an Anglican bishop and no enemy of his own office, called the identity of the two titles in the New Testament a fact generally recognized by theologians of every shade of opinion.

The earliest witnesses divide. Clement of Rome, writing to Corinth about the year 96, knows bishops and deacons appointed by the apostles, calls the men who hold the bishop's office presbyters in the same breath, and addresses a Corinth governed by its presbyters, with no single bishop anywhere in the letter. The Didache tells congregations to appoint for themselves bishops and deacons. St. Polycarp writes to Philippi through its presbyters and deacons, names no bishop there, and does not call himself one. Against this stands St. Ignatius of Antioch, traditionally about the year 107, though some recent scholarship dates the letters a generation later, who assumes a single bishop in every Asian church he writes to, Ephesus, Magnesia, Tralles, Smyrna, and urges that nothing be done apart from him. Two details qualify the witness. His letter to Rome, alone of the seven, mentions no bishop at all. And the imagery of his strongest passage runs the other way from the later theory, for in Magnesians the bishop presides after the likeness of God and the presbyters after the likeness of the council of the apostles, so that the imagery of succession from the apostolic college attaches, in the episcopate's greatest ancient witness, to the presbyterate.

The doctors then say plainly what the record shows. St. Jerome, in the letter to Evangelus, argues from the texts above that the presbyter is the same as the bishop, and that the bishop rose out of the presbyterate as a remedy for schism, the presbyters of Alexandria from St. Mark until Heraclas and Dionysius choosing one of their own number and naming him bishop, as an army makes an emperor. St. John Chrysostom teaches that the presbyters were of old called bishops and the bishops presbyters. And the reading did not die with the Fathers, for the schoolmen debated whether the episcopate is a distinct order at all or the perfection of the one priesthood, and Aquinas in the Supplement took the second view. The first reading's honest limits are recorded with its strengths. St. Jerome's Alexandria is election, not consecration, his verbs are chose and named, and the reports of presbyters laying on hands there rest on later and disputed sources. And St. Jerome himself asks what a bishop does that a presbyter does not, excepting ordination, reserving in that clause the very act in question.

The practice record is thin but real. Cassian tells of the desert presbyter Paphnutius who preferred Daniel to the diaconate and promoted him to the priesthood, no bishop appearing in the account, though the verbs can be read as nomination rather than ordination and have been read both ways. Harder to explain away are the popes. In 1400 Boniface IX granted an abbot who was no bishop the power to confer on his canons all orders up to and including the priesthood, and when the grant was revoked three years later on jurisdictional complaint, the ordinations already given were not declared invalid. In 1427 Martin V gave a Cistercian abbot the same power for five years. In 1489 Innocent VIII licensed the abbot of C\^iteaux and the four proto-abbots to ordain their monks subdeacons and deacons, a privilege used into the seventeenth century and formally ended only in 1902. These acts stand in the registers and in Denzinger (DH 1145 to 1146, 1290, 1435), and whatever theology explains them, a doctrine on which presbyteral ordination is simply impossible does not.

On these facts one of us reads the episcopate as St. Jerome read it, the ruling presbyterate, an arrangement of the Church's own wisdom for unity, so that the presbyteral successions of the Reformation are real successions bearing a real defect of order rather than nothing at all. Another of us reads the same facts as John Henry Newman would, the seed of oversight in the apostles and in Timothy and Titus, grown so fast and so universally that within a lifetime of the apostles St. Ignatius can assume it across Asia, which does not happen to innovations, so the episcopate is an authentic development received by the whole Body and kept among this paper's marks. And the reception case is strong, for whatever the first century was, the threefold order was universal by the end of the second and unanimous until the sixteenth, a reception as deep as any this paper grades. The first reading's reply is that the wobble sits exactly at the seed, where this paper's own rules look first, and that what the whole Body received may have been received as the Church's right ordering rather than as the being of the ministry itself, order and discipline rather than deposit, a distinction the paper already uses when it grades personal anathemas. And the reply carries its own guard, lest it be turned against section 2. What licenses the first reading is counter-witness at the seed, St. Jerome and St. John Chrysostom themselves reading the titles as one, Clement and the Didache and St. Polycarp beside them. On the reservation of orders to men there is no such witness at all, no Father reads the offices as open, and the deaconess canons show the boundary drawn on purpose. A question wobbles at the seed only when the seed's own doctors wobble, and here they do not.

Pressed as far as it will go, the first reading changes the map without abolishing it. The episcopate remains the received shape of the Church and its recovery is owed everywhere, which is the old Anglican position, Hooker's, that the government of bishops is the Church's well-being while presbyteral orders are not therefore null. The five marks then read the episcopate as owed rather than as the boundary of being, the first tier widens to the confessional bodies of presbyteral succession that pass the moral condition and the remaining marks, the receiver gains members, and the mutual act of section 7 grows arms it does not now have. And the two readings are not two roads. They are one road walked at two widths, both defensible and both respectable. The paper's body walks the narrower width and keeps its focus, the marks and the tiers standing as written on the second reading, while this afterword records the wider walking at full length because it is real, its fathers are St. Jerome and St. John Chrysostom, its registers are papal, and honest men inside this confession weigh it differently and remain brothers.

One line remains, and the wider reading draws it as firmly as the narrower, so it is drawn here where the tier stands widest. The widening reaches the confessional bodies, the churches of the Augsburg Confession and the old Reformed standards, because what they kept is the question's very subject, an ordered ministry conferred by ordination in succession, the creed confessed, the font given to infants as the whole Church gave it, the Supper a means of grace, the moral law unbroken, though the remaining marks still sort within that reach, for the Supper of the Reformed standards is a means of grace but not the presence Augsburg confesses. It does not reach the fundamentalism and revivalist evangelicalism of the last two centuries, and the reason is not taste but reception. The test of this whole paper is what the whole Body received in its first thousand years, and set beside it the revivalist's faith and practice appear quite different, not in fervor but in shape. No creed recited, no font for the covenant's children, no Supper beyond a memorial, no ordered ministry beyond a call answered and a pulpit founded. The confessional Lutheran stands near the first tier because he holds most of what the undivided Church held. The revivalist holds a faith the first millennium would not recognize as its own, sincere, often fervent, sometimes holier than its critics, but not the received thing, and a paper whose one criterion is reception cannot pretend otherwise. The line is not contempt. It is the same line the paper draws everywhere, between what the whole received and what a part devised, and it cuts against the newest devisings exactly as it cuts against Rome's.

Like every question of its size the wider reading belongs to the reception of the whole Body, and it is submitted there. The Lutheran who holds everything here save the government of bishops should know that the ground he stands on is not nothing. St. Jerome stood near it.
